
注意，若矩阵 $\Ub$ 左标准正交，则 $\Ub\Ub^\ts$ 未必是单位矩阵，即 
$$
\begin{tikzpicture}[baseline=\baseline]
    \tikzset{tensor/.style = {minimum size = 0.5cm,shape = circle,thick,draw=black,inner sep = 0pt}, edge/.style   = {thick,line width=.4mm},every loop/.style={}}
    \def\y{0}
    \def\x{0};

    \node[tensor, path picture={
    \fill[fill=red!30!white](path picture bounding box.south east)
    -- 
    (path picture bounding box.north west) --
    (path picture bounding box.north east) -- cycle;
    }, shift={(\x,\y)}] (U) {$\Ub$};
    
    \node[tensor, path picture={
    \fill[fill=red!30!white](path picture bounding box.south west)
        -- 
    (path picture bounding box.north east) --
    (path picture bounding box.north west) -- cycle;
    }, shift={(\x+1,\y)}] (Ut) {$\Ub$};

    \draw[edge] (U) -- (Ut) node [midway,above] {\scalebox{0.7}{\dimcolor{$n$}}};
    \draw[edge](Ut) -- (\x+2,\y) node [midway,above] {\scalebox{0.7}{\dimcolor{$m$}}};
    \draw[edge](U) -- (\x-1,\y) node [midway,above] {\scalebox{0.7}{\dimcolor{$m$}}};
    \node[draw=none] () at (\x+4.3,\y) {\scalebox{0.8}{\dimcolor{$=$}}};
    \node[draw=none] () at (\x+2.3,\y) {\scalebox{0.8}{\dimcolor{$\neq$}}};
    \draw[edge](\x+2.5,\y) -- (\x+4,\y)node [midway,above] {\scalebox{0.7}{\dimcolor{$m$}}};
\end{tikzpicture}\Ib_m,$$ 只有当 $m=n$ 时才会如此。 



借助这些视觉线索，QR 分解与 SVD 对应的张量网络如下所示：
\begin{itemize}
    \item\textbf{QR 分解} 给定 $\Ab\in\RR^{d_1\times d_2}$，其~（秩揭示）QR 分解为 
    
$$
\begin{tikzpicture}[baseline=\baseline]
    \tikzset{tensor/.style = {minimum size = 0.5cm,shape = circle,thick,draw=black,inner sep = 0pt}, edge/.style   = {thick,line width=.4mm},every loop/.style={}}
    \def\y{-0.2}
    \def\x{0}
    \node[tensor,fill=red!30!white] (A) at (\x,\y+0.2){$\Ab$};
    \draw[edge] (A) -- (\x-0.8,\y+0.2)  node [midway,above] {\scalebox{0.7}{\dimcolor{$d_1$}}};
    \draw[edge] (A) -- (\x+0.8,\y+0.2)  node [midway,above] {\scalebox{0.7}{\dimcolor{$d_2$}}};
    \node[draw = none] at (\x+1.2,\y+0.2){$=$};
    \node[tensor, path picture={
    \fill[fill=red!30!white](path picture bounding box.south east)
    -- 
    (path picture bounding box.north west) --
    (path picture bounding box.north east) -- cycle;
    }, shift={(\x+2.5,\y+0.2)}] (Q) {$\Qb$};
    \node[tensor,fill=blue!30!white] (R) at (\x+3.5,\y+0.2){$\Rb$};
    \draw[edge] (Q) -- (R)  node [midway,above] {\scalebox{0.7}{\dimcolor{$R$}}};
    \draw[edge] (Q) -- (\x+1.7,\y+0.2)  node [midway,above] {\scalebox{0.7}{\dimcolor{$d_1$}}};
    \draw[edge] (R) -- (\x+4.3,\y+0.2)  node [midway,above] {\scalebox{0.7}{\dimcolor{$d_2$}}};
\end{tikzpicture},$$
其中 
$\begin{tikzpicture}[baseline=\baseline]
    \tikzset{tensor/.style = {minimum size = 0.5cm,shape = circle,thick,draw=black,inner sep = 0pt}, edge/.style   = {thick,line width=.4mm},every loop/.style={}}
    \def\y{0}
    \def\x{0}
    \node[tensor, path picture={
    \fill[fill=red!30!white](path picture bounding box.south east)
    -- 
    (path picture bounding box.north west) --
    (path picture bounding box.north east) -- cycle;
    }, shift={(\x,\y)}] (Q) {$\Qb$};
    \draw[edge] (Q) -- (\x+0.7,\y)  node [midway,above] {\scalebox{0.7}{\dimcolor{$R$}}};
    \draw[edge] (Q) -- (\x-0.7,\y)  node [midway,above] {\scalebox{0.7}{\dimcolor{$d_1$}}};
\end{tikzpicture}
\in\RR^{d_1\times R}$
是左标准正交矩阵，即 $\Qb^\ts\Qb = \Ib_R$；$\Rb\in\RR^{R\times d_2}$ 为上三角矩阵~（这一点在张量网络表示中并不明显），$R$ 是 $\Ab$ 的秩。 
\item \textbf{奇异值分解~（SVD）} 给定 $\Ab\in\RR^{d_1\times d_2}$，其~（紧凑）SVD 为 
\begin{align*}
\begin{tikzpicture}[baseline=\baseline]
    \tikzset{tensor/.style = {minimum size = 0.5cm,shape = circle,thick,draw=black,inner sep = 0pt}, edge/.style   = {thick,line width=.4mm},every loop/.style={}}
    \def\y{0}
    \def\x{0}
    \node[tensor,fill=red!30!white] (A) at (\x,\y){$\Ab$};
    \draw[edge] (A) -- (\x-0.8,\y)  node [midway,above] {\scalebox{0.7}{\dimcolor{$d_1$}}};
    \draw[edge] (A) -- (\x+0.8,\y)  node [midway,above] {\scalebox{0.7}{\dimcolor{$d_2$}}};
    \node[draw = none] at (\x+1.35,\y){$=$};
    \node[tensor, path picture={
    \fill[fill=red!30!white](path picture bounding box.south east)
    -- 
    (path picture bounding box.north west) --
    (path picture bounding box.north east) -- cycle;
    }, shift={(\x+3,\y)}] (U) {$\Ub$};
    \node[tensor,fill=blue!20!red!40!white] (sigma) at (\x+4.2,\y){\scalebox{0.85}{}};
    \node[tensor, path picture={
    \fill[fill=red!30!white](path picture bounding box.south west)
        -- 
    (path picture bounding box.north west) --
    (path picture bounding box.north east) -- cycle;
    }, shift={(\x+5.2,\y)}] (V) {$\Vb$};
    \draw[edge=0.5] (sigma.south west) -- (sigma.north        east);
    \draw[edge] (U) -- (sigma) node [midway,above] {\scalebox{0.7}{\dimcolor{$R$}}};
    \draw[edge] (sigma) -- (V)node [midway,above] {\scalebox{0.7}{\dimcolor{$R$}}};
    \draw[edge](U) -- (\x+2,\y) node [midway,above] {\scalebox{0.7}{\dimcolor{$d_1$}}};
    \draw[edge](V) -- (\x+6.2,\y) node [midway,above] {\scalebox{0.7}{\dimcolor{$d_2$}}};
\end{tikzpicture}~,
\end{align*}
其中 $\Ub\in\RR^{d_1\times R}$ 左标准正交~（$\Ub^\ts\Ub= \Ib_R$），$\Vb\in\RR^{R\times d_2}$ 右标准正交~（$\Vb\Vb^\ts = \Ib_{R}$），$\Sigmab\in\RR^{R\times R}
=
\begin{tikzpicture}[baseline=\baseline]
    \tikzset{tensor/.style = {minimum size = 0.5cm,shape = circle,thick,draw=black,inner sep = 0pt}, edge/.style   = {thick,line width=.4mm},every loop/.style={}}
    \def\y{0}
    \def\x{0}
    \node[tensor,fill=blue!20!red!40!white] (sigma) at (\x,\y){\scalebox{0.85}{}};
    \draw[edge=0.5] (sigma.south west) -- (sigma.north        east);
    \draw[edge](sigma) -- (\x-0.7,\y);
    \draw[edge](sigma) --(\x+0.7,\y);
\end{tikzpicture}$ 为对角矩阵，$R$ 是 $\Ab$ 的秩。 
\end{itemize}

